\documentclass[11pt]{article}
\usepackage{fontspec}
\directlua{luaotfload.add_fallback("hebfb",{"DejaVu Sans:mode=harf"})}
\usepackage{polyglossia}
\setdefaultlanguage{hebrew}
\setotherlanguage{english}
\setmainfont{Noto Sans Hebrew}[Script=Hebrew,RawFeature={fallback=hebfb}]
\newfontfamily\codefont{DejaVu Sans Mono}[RawFeature={fallback=hebfb}]
\usepackage[a4paper,margin=18mm]{geometry}
\usepackage[table]{xcolor}
\usepackage{mdframed}
\usepackage{tabularx}
\usepackage{xltabular}
\usepackage{array}
\usepackage{enumitem}
\usepackage{fancyhdr}
\setlist{leftmargin=1.6em,itemsep=1pt,topsep=2pt}
% colors
\definecolor{h1}{HTML}{1E3A8A}\definecolor{h2}{HTML}{1E40AF}\definecolor{h3}{HTML}{0F172A}
\definecolor{subgray}{HTML}{64748B}\definecolor{hdrbg}{HTML}{E0E7FF}
\definecolor{cfg}{HTML}{E2E8F0}\definecolor{cbg}{HTML}{0F172A}\definecolor{cbold}{HTML}{7DD3FC}\definecolor{ccom}{HTML}{94A3B8}
\definecolor{metabg}{HTML}{F1F5F9}
\newcommand{\enLR}[1]{\textenglish{#1}}
\newcommand{\code}[1]{\textenglish{\codefont #1}}
\newcommand{\chapterheading}[1]{\vspace{2pt}{\LARGE\bfseries\color{h1}#1\par}\vskip2pt{\color{h1}\hrule height 2pt}\vskip6pt}
\newcommand{\sectionheading}[1]{\vskip10pt\penalty-200{\large\bfseries\color{h2}#1\par}\vskip3pt}
\newcommand{\subheading}[1]{\vskip6pt{\bfseries\color{h3}#1\par}\vskip2pt}
\newcommand{\boxtitle}[1]{{\bfseries\color{boxti}#1}}
% box environments (right border, tinted bg)
\newmdenv[topline=false,bottomline=false,leftline=false,rightline=true,linewidth=2pt,innermargin=0pt,skipabove=6pt,skipbelow=6pt,innertopmargin=6pt,innerbottommargin=6pt,innerleftmargin=8pt,innerrightmargin=8pt]{genericbox}
\definecolor{faqbg}{HTML}{ECFDF5}\definecolor{faqfr}{HTML}{10B981}\definecolor{faqti}{HTML}{065F46}
\definecolor{misbg}{HTML}{FEF2F2}\definecolor{misfr}{HTML}{EF4444}\definecolor{misti}{HTML}{991B1B}
\definecolor{tipbg}{HTML}{FFFBEB}\definecolor{tipfr}{HTML}{F59E0B}\definecolor{tipti}{HTML}{92400E}
\definecolor{exabg}{HTML}{EFF6FF}\definecolor{exafr}{HTML}{3B82F6}\definecolor{exati}{HTML}{1E3A8A}
\definecolor{defbg}{HTML}{F5F3FF}\definecolor{deffr}{HTML}{8B5CF6}\definecolor{defti}{HTML}{5B21B6}
\definecolor{stobg}{HTML}{FDF4FF}\definecolor{stofr}{HTML}{D946EF}\definecolor{stoti}{HTML}{86198F}
\definecolor{exbg}{HTML}{F0FDFA}\definecolor{exfr}{HTML}{14B8A6}\definecolor{exti}{HTML}{115E59}
\newenvironment{faqbox}{\colorlet{boxti}{faqti}\begin{mdframed}[backgroundcolor=faqbg,linecolor=faqfr,rightline=true,leftline=false,topline=false,bottomline=false,linewidth=2pt,innertopmargin=6pt,innerbottommargin=6pt,innerleftmargin=8pt,innerrightmargin=8pt,skipabove=6pt,skipbelow=6pt]}{\end{mdframed}}
\newenvironment{mistakebox}{\colorlet{boxti}{misti}\begin{mdframed}[backgroundcolor=misbg,linecolor=misfr,rightline=true,leftline=false,topline=false,bottomline=false,linewidth=2pt,innertopmargin=6pt,innerbottommargin=6pt,innerleftmargin=8pt,innerrightmargin=8pt,skipabove=6pt,skipbelow=6pt]}{\end{mdframed}}
\newenvironment{tipbox}{\colorlet{boxti}{tipti}\begin{mdframed}[backgroundcolor=tipbg,linecolor=tipfr,rightline=true,leftline=false,topline=false,bottomline=false,linewidth=2pt,innertopmargin=6pt,innerbottommargin=6pt,innerleftmargin=8pt,innerrightmargin=8pt,skipabove=6pt,skipbelow=6pt]}{\end{mdframed}}
\newenvironment{exambox}{\colorlet{boxti}{exati}\begin{mdframed}[backgroundcolor=exabg,linecolor=exafr,rightline=true,leftline=false,topline=false,bottomline=false,linewidth=2pt,innertopmargin=6pt,innerbottommargin=6pt,innerleftmargin=8pt,innerrightmargin=8pt,skipabove=6pt,skipbelow=6pt]}{\end{mdframed}}
\newenvironment{defbox}{\colorlet{boxti}{defti}\begin{mdframed}[backgroundcolor=defbg,linecolor=deffr,rightline=true,leftline=false,topline=false,bottomline=false,linewidth=2pt,innertopmargin=6pt,innerbottommargin=6pt,innerleftmargin=8pt,innerrightmargin=8pt,skipabove=6pt,skipbelow=6pt]}{\end{mdframed}}
\newenvironment{storybox}{\colorlet{boxti}{stoti}\begin{mdframed}[backgroundcolor=stobg,linecolor=stofr,rightline=true,leftline=false,topline=false,bottomline=false,linewidth=2pt,innertopmargin=6pt,innerbottommargin=6pt,innerleftmargin=8pt,innerrightmargin=8pt,skipabove=6pt,skipbelow=6pt]}{\end{mdframed}}
\newenvironment{exbox}{\colorlet{boxti}{exti}\begin{mdframed}[backgroundcolor=exbg,linecolor=exfr,rightline=true,leftline=false,topline=false,bottomline=false,linewidth=2pt,innertopmargin=6pt,innerbottommargin=6pt,innerleftmargin=8pt,innerrightmargin=8pt,skipabove=6pt,skipbelow=6pt]}{\end{mdframed}}
\newenvironment{metabox}{\begin{mdframed}[backgroundcolor=metabg,linewidth=0pt,innertopmargin=5pt,innerbottommargin=5pt,innerleftmargin=8pt,innerrightmargin=8pt,skipabove=4pt,skipbelow=8pt]}{\end{mdframed}}
\newmdenv[backgroundcolor=cbg,linewidth=0pt,innertopmargin=6pt,innerbottommargin=6pt,innerleftmargin=8pt,innerrightmargin=8pt,skipabove=6pt,skipbelow=8pt]{codebox}
\setlength{\parindent}{0pt}\setlength{\parskip}{4pt}
\renewcommand{\thepage}{\enLR{\arabic{page}}}
\pagestyle{fancy}\fancyhf{}\fancyfoot[C]{\thepage}\renewcommand{\headrulewidth}{0pt}
\begin{document}
\sloppy
\chapterheading{פרק \enLR{3} – מודל ה-\enLR{AAA}}{\small\color{subgray}\enLR{11} שעות עיוני + \enLR{3} מעשי · שבועות \enLR{7}–\enLR{9}\par}\vskip4pt

\begin{defbox}\boxtitle{איך לקרוא את הפרק הזה}\par  בפרק הקודם הקשחנו נתב אחד ונתנו לו סיסמה מקומית. אבל מה עושים כשיש עשרות נתבים ומתגים ועשרות אנשים שצריכים גישה? כאן נכנס מודל ה-\enLR{AAA.} הפרק כתוב לקריאה רציפה: נתחיל מהבעיה שהמודל בא לפתור, נבין את שלוש ה-\enLR{A} במילים פשוטות, ורק אז נגיע לפרוטוקולים ולפקודות. כל מונח מוסבר בפעם הראשונה שהוא מופיע.\end{defbox}

\sectionheading{\enLR{3.1} הבעיה: איך מנהלים גישה בעשרות מכשירים?}

נזכיר מהפרק הקודם: הגדרנו על הנתב \textbf{משתמש מקומי \enLR{(local user)}} – שם וסיסמה ששמורים בתוך הנתב עצמו. זה עובד מצוין לנתב אחד. אבל דמיינו רשת של בית ספר או חברה עם \enLR{50} נתבים ומתגים, ו-\enLR{10} אנשי צוות שצריכים לנהל אותם. אם כל מכשיר מחזיק את רשימת המשתמשים שלו, נוצרות שתי בעיות כואבות: ראשית, כדי להוסיף עובד חדש צריך להיכנס ידנית ל-\enLR{50} מכשירים; שנית – וזה החמור – כשעובד \underline{עוזב}, צריך לזכור למחוק אותו מכל \enLR{50} המכשירים, ואם שוכחים אפילו אחד, נשארת דלת פתוחה.\par

הפתרון: \textbf{ניהול מרכזי}. במקום שכל מכשיר יחזיק רשימת משתמשים משלו, יש \underline{מקום אחד} – שרת מרכזי – שכל המכשירים שואלים אותו "האם להכניס את המשתמש הזה, ומה מותר לו?". הוספה או מחיקה של עובד נעשית פעם אחת, במקום אחד, ומשפיעה על כל הרשת. זו בדיוק המסגרת שנקראת \textbf{\enLR{AAA}}.\par

\sectionheading{\enLR{3.2} שלוש ה-\enLR{A}: מה קורה כשמתחברים}

כשאתם נכנסים לחשבון כלשהו – מייל, בנק, נתב – קורים שלושה דברים נפרדים שאנשים נוטים לבלבל ביניהם. מודל \enLR{AAA} (שלושה ראשי תיבות באנגלית) בדיוק מסדר אותם:\par

\begin{itemize}
\item \textbf{אימות \enLR{(Authentication)}} – עונה על השאלה \underline{"מי אתה?"}. זהו שלב הוכחת הזהות, למשל בעזרת שם משתמש וסיסמה.
\item \textbf{הרשאה \enLR{(Authorization)}} – עונה על השאלה \underline{"מה מותר לך לעשות?"}. אחרי שידוע מי אתה, ההרשאה קובעת מה מותר – האם רק לצפות, או גם לשנות הכול? לאילו מכשירים יש לך גישה?
\item \textbf{רישום \enLR{(Accounting)}} – עונה על השאלה \underline{"מה עשית, ומתי?"}. זהו תיעוד של פעולות המשתמש: מתי נכנס, אילו פקודות הקליד, מתי יצא.
\end{itemize}

\begin{tipbox}\boxtitle{המשל שעוזר לזכור: שדה התעופה}\par  דמיינו כניסה לטיסה. \textbf{אימות} = בדיקת הדרכון בכניסה – מוודאים שאתם מי שאתם טוענים שאתם. \textbf{הרשאה} = כרטיס הטיסה – גם אחרי שידוע מי אתם, מותר לכם לעלות רק לטיסה שלכם, לא לכל מטוס בשדה. \textbf{רישום} = רשימת הנוסעים שהמדינה שומרת – מי עלה, מתי, עם כמה מזוודות. שימו לב לסדר ההגיוני: אי אפשר לתת הרשאה בלי לדעת קודם מי האדם (אימות), ואי אפשר לרשום פעולות של מישהו אלמוני.\end{tipbox}

\subheading{גורמי אימות – שלוש דרכים להוכיח מי אתה}

איך מוכיחים זהות? יש שלושה "גורמים" \enLR{(factors)}, ולכל אחד יתרון וחיסרון:\par

\begin{itemize}
\item \textbf{משהו שאתה יודע} – סיסמה, קוד \enLR{PIN.} קל לשימוש, אבל אפשר לגנוב או לנחש.
\item \textbf{משהו שיש לך} – הטלפון (שמקבל קוד חד-פעמי), כרטיס חכם, טוקן. קשה יותר לגנוב מרחוק.
\item \textbf{משהו שאתה} – ביומטריה: טביעת אצבע, זיהוי פנים. קשה מאוד לזייף.
\end{itemize}

כשמשלבים \underline{שני גורמים מסוגים שונים} מקבלים \textbf{אימות רב-שלבי \enLR{(MFA)}}, או במקרה של שניים בדיוק – \textbf{\enLR{2FA}}. זו הגנה חזקה מאוד: גם אם תוקף גנב את הסיסמה שלכם (משהו שאתם יודעים), עדיין חסר לו הטלפון (משהו שיש לכם). שימו לב: שתי סיסמאות אינן \enLR{MFA} – חייבים גורמים משני \emph{סוגים} שונים.\par

\sectionheading{\enLR{3.3} שני מצבים: \enLR{AAA} מקומי מול מבוסס-שרת}

יש שתי דרכים להפעיל \enLR{AAA.} ב-\textbf{\enLR{AAA} מקומי \enLR{(Local)}} המכשיר מאמת מול רשימת משתמשים ששמורה בו עצמו – פשוט, מתאים לרשת קטנה, וגם משמש כ"גיבוי" למקרה שהשרת המרכזי לא זמין. ב-\textbf{\enLR{AAA} מבוסס-שרת \enLR{(Server-based)}} המכשיר פונה ל\underline{שרת מרכזי} ושואל אותו – זו הגישה הנכונה לארגון, כי היא נותנת ניהול אחיד ורישום מרוכז.\par

כדי לדבר על זה נכון צריך שני מושגים: המכשיר שפונה לשרת (הנתב או המתג) נקרא \textbf{\enLR{NAS (Network Access Server)}} או פשוט "לקוח ה-\enLR{AAA".} והחיבור בין הנתב לשרת מאובטח בעזרת \textbf{מפתח משותף \enLR{(shared key / secret)}} – סוד שמוגדר \underline{גם בנתב וגם בשרת}, ומשמש להצפין את התקשורת ביניהם ולהוכיח לשרת שהנתב הוא לקוח מורשה שלו.\par

\begin{mistakebox}\boxtitle{בלבול נפוץ: המפתח אינו הסיסמה של המשתמש!}\par  המפתח המשותף \enLR{(key)} הוא סוד בין \underline{הנתב לשרת}. הסיסמאות של המשתמשים שמורות בשרת. אלה שני דברים שונים לגמרי: המפתח מאמת את הציוד, הסיסמה מאמתת את האדם. בתכנית הלימודים הביטוי "מפתחות מוצפנים" מתייחס למפתח המשותף הזה.\end{mistakebox}

\begin{storybox}\boxtitle{מהעולם האמיתי: מנהל הרשת שפוטר ביום שלישי}\par  בחברת שירותי אינטרנט פיטרו מנהל רשת. המנכ"ל ביקש מה-\enLR{IT "}לבטל לו את הגישה", וה-\enLR{IT} מחק את חשבון ה-\enLR{Windows} שלו. אבל ב-\enLR{40} הנתבים והמתגים היו חשבונות \underline{מקומיים} שהמנהל עצמו יצר. שלושה ימים אחר כך, ב-\enLR{2} בלילה, כל הנתבים "איבדו" את ההגדרות שלהם – ואף אחד לא ידע מי עשה זאת, כי לא היה רישום \enLR{(Accounting)} וכל החשבונות נראו כמו "\enLR{admin".} עם \enLR{AAA} מבוסס-שרת, לחיצה אחת הייתה מנתקת אותו מכל הציוד, והרישום היה מראה שם מלא, שעה ופקודה. זה בדיוק ה"למה" של הפרק.\end{storybox}

\sectionheading{\enLR{3.4 RADIUS} מול \enLR{TACACS}+ – שני הפרוטוקולים}

כשהנתב פונה לשרת המרכזי, הוא מדבר איתו באחד משני פרוטוקולים (שפות). שניהם עושים \enLR{AAA}, אבל יש ביניהם הבדלים חשובים שנשאלים כמעט תמיד בבגרות. נכיר כל אחד, ואז נשווה.\par

\textbf{\enLR{RADIUS}} הוא תקן פתוח (לא שייך לחברה מסוימת). הוא רץ על פרוטוקול התעבורה \textbf{\enLR{UDP}} (פרוטוקול מהיר וחסר-חיבור), ומצפין \underline{רק את שדה הסיסמה} בבקשה – שאר הפרטים (כולל שם המשתמש) עוברים גלויים. \enLR{RADIUS} גם \underline{משלב} את האימות וההרשאה יחד: כשהשרת מאשר את המשתמש, התשובה כוללת כבר גם את ההרשאות. השימוש האופייני שלו: \textbf{גישת משתמשים לרשת} – למשל אימות ל-\enLR{Wi-Fi} ארגוני \enLR{(802.1X)} או ל-\enLR{VPN} מרחוק.\par

\textbf{\enLR{TACACS}+} הוא פרוטוקול של חברת סיסקו. הוא רץ על \textbf{\enLR{TCP}} (פרוטוקול אמין שמנהל "שיחה"), פורט \enLR{49}, ומצפין את \underline{כל} גוף החבילה – לא רק הסיסמה. הוא גם \underline{מפריד} בין שלושת השירותים, כך שאפשר, למשל, לאמת מול מקור אחד ולאשר כל פקודה בנפרד מול השרת. השימוש האופייני שלו: \textbf{ניהול ציוד רשת} – מנהלים שמתחברים לנתבים ולמתגים, עם אפשרות לאשר כל פקודה שהם מקלידים.\par

\begin{xltabular}{\linewidth}{|>{\setRTL\raggedleft\arraybackslash}X|>{\setRTL\raggedleft\arraybackslash}X|>{\setRTL\raggedleft\arraybackslash}X|}
\hline
\rowcolor{hdrbg}\textbf{\enLR{TACACS}+} & \textbf{\enLR{RADIUS}} & \textbf{} \\
\hline
\endhead
סיסקו & תקן פתוח \enLR{(IETF)} & מי פיתח \\
\hline
\enLR{TCP} (פורט \enLR{49)} & \enLR{UDP (1812/1813)} & תעבורה \\
\hline
כל גוף החבילה & רק שדה הסיסמה & הצפנה \\
\hline
שלושתם נפרדים & אימות והרשאה משולבים & הפרדת השירותים \\
\hline
נתמך & לא נתמך & אישור לכל פקודה \\
\hline
ניהול ציוד & גישת משתמשים \enLR{(Wi-Fi/VPN)} & שימוש אופייני \\
\hline
\end{xltabular}

\begin{tipbox}\boxtitle{למה \enLR{RADIUS} על \enLR{UDP} אם \enLR{TCP} אמין יותר?}\par  שאלה טובה שתלמידים שואלים. \enLR{RADIUS} תוכנן ב-\enLR{1991} לשרתי דיאל-אפ עם אלפי חיבורים בו-זמנית; \enLR{UDP} קל ומהיר, והלקוח פשוט שולח שוב אם לא קיבל תשובה. \enLR{TACACS}+ בחר \enLR{TCP} כי הוא מנהל שיחה רב-שלבית (שאלה–תשובה–שאלה) ורוצה לדעת מיד אם השרת נפל. שניהם עובדים מצוין היום; ההבדל המעשי החשוב הוא ההצפנה המלאה והפרדת השירותים של \enLR{TACACS+.}\end{tipbox}

\sectionheading{\enLR{3.5} השרתים בפועל: \enLR{ACS} ו-\enLR{ISE}}

מי מריץ את שירות ה-\enLR{AAA}? אצל סיסקו, השרת ההיסטורי נקרא \textbf{\enLR{ACS (Access Control Server)}}, והיורש המודרני שלו נקרא \textbf{\enLR{ISE (Identity Services Engine)}} – שרת מדיניות שעושה \enLR{AAA}, מנהל \enLR{802.1X}, ואפילו בודק את מצב המכשירים שמתחברים. יש גם חלופות שאינן של סיסקו: \textbf{\enLR{FreeRADIUS}} (קוד פתוח) ו-\textbf{\enLR{Windows NPS}} (שרת \enLR{RADIUS} שמובנה ב-\enLR{Windows Server).} היתרון הגדול: השרת יכול לאמת מול מאגר המשתמשים הקיים של הארגון \enLR{(Active Directory)}, כך שהמורים משתמשים באותו שם וסיסמה בכל מקום.\par

\sectionheading{\enLR{3.6} הגדרה בפועל – שלב אחר שלב}

עכשיו שהבנו את הרעיון, נראה איך מגדירים. הצעד הראשון תמיד: להפעיל את מנגנון ה-\enLR{AAA} בפקודה \code{aaa new-model}. מרגע זה, כל קווי הגישה מתחילים להשתמש ב-\enLR{AAA} במקום בשיטה הישנה.\par

\begin{codebox}\begin{english}\codefont\footnotesize\color{cfg}\setlength{\parindent}{0pt}
R1(config)\# username admin secret P@ss        \textcolor{ccom}{\texthebrew{! משתמש מקומי – גיבוי, קודם כול!}}\\
R1(config)\# \textcolor{cbold}{\textbf{aaa new-model}}                    \textcolor{ccom}{\texthebrew{! מפעיל את מסגרת ה-\enLR{AAA}}}\\
R1(config)\# tacacs-server host 10.0.0.5 key Key123  \textcolor{ccom}{\texthebrew{! שרת \enLR{TACACS}+ והמפתח המשותף}}\\
R1(config)\# \textcolor{cbold}{\textbf{aaa authentication login default group tacacs+ local}}
\end{english}\end{codebox}

השורה האחרונה היא "רשימת שיטות" \enLR{(method list)}: היא אומרת "כדי לאמת התחברות, נסה קודם את שרת ה-\enLR{TACACS}+; ואם הוא \underline{לא זמין} – השתמש במשתמש המקומי". הגיבוי המקומי (\code{local}) קריטי: בלעדיו, אם השרת נופל, אף אחד לא יוכל להיכנס.\par

\begin{mistakebox}\boxtitle{הטעות המסוכנת ביותר בפרק – היזהרו!}\par  הפקודה \code{aaa new-model} משנה מיד את התנהגות כל הקווים. אם תקלידו אותה \underline{בלי} להגדיר קודם משתמש מקומי ובלי גיבוי \code{local} ברשימת השיטות, ואז תתנתקו – \textbf{תינעלו מחוץ לנתב}. כלל ברזל: קודם הגדירו \code{username}, אחר כך \code{aaa new-model}, תמיד סיימו את רשימת השיטות ב-\code{local}, והשאירו חלון חיבור פתוח עד שווידאתם בחלון שני שהכול עובד.\end{mistakebox}

נקודה חשובה שנשאלת בבגרות: מעבר לשיטה הבאה ברשימה קורה \underline{רק} כאשר השיטה הקודמת \textbf{אינה זמינה} (השרת לא עונה). אם השרת עונה "סיסמה שגויה" – המשתמש נדחה, ואין מעבר לשיטה הבאה. כלומר \code{local} הוא גיבוי לתקלת-שרת, לא "ניסיון שני" אחרי סיסמה שגויה.\par

הגדרת שרת \enLR{RADIUS} דומה מאוד, רק עם מילת המפתח \code{radius}:\par

\begin{codebox}\begin{english}\codefont\footnotesize\color{cfg}\setlength{\parindent}{0pt}
R1(config)\# radius-server host 10.0.0.6 key R@diusK3y\\
R1(config)\# aaa authentication login default group radius local
\end{english}\end{codebox}

\subheading{הרשאה ורישום}

עד כה עסקנו באימות (ה-\enLR{A} הראשון). מוסיפים את שני האחרים בשתי שורות דומות. \textbf{הרשאה \enLR{(Authorization)}} קובעת מה מותר למשתמש לעשות אחרי שנכנס – למשל, האם הוא מקבל גישת מנהל (רמה \enLR{15)} ישירות:\par

\begin{codebox}\begin{english}\codefont\footnotesize\color{cfg}\setlength{\parindent}{0pt}
R1(config)\# aaa authorization exec default group tacacs+ local
\end{english}\end{codebox}

\textbf{רישום \enLR{(Accounting)}} מתעד את הפעולות. למשל, לתעד כל פקודה ברמה \enLR{15} (מנהל) – כדי שתמיד נדע "מי הקליד \enLR{reload}?":\par

\begin{codebox}\begin{english}\codefont\footnotesize\color{cfg}\setlength{\parindent}{0pt}
R1(config)\# aaa accounting commands 15 default start-stop group tacacs+
\end{english}\end{codebox}

\sectionheading{\enLR{3.7 802.1X} – \enLR{AAA} שמגן על שקע ברשת}

\enLR{RADIUS} משמש גם לדבר חשוב נוסף: לאמת \underline{מכשיר} לפני שהוא בכלל מקבל גישה לרשת. התקן שעושה זאת נקרא \textbf{\enLR{802.1X}}, ויש בו שלושה תפקידים: ה-\textbf{\enLR{Supplicant}} (המחשב או הטלפון שרוצה להתחבר), ה-\textbf{\enLR{Authenticator}} (המתג או נקודת הגישה, ש\underline{חוסם את הפורט} עד שהאימות מצליח), וה-\textbf{שרת האימות} \enLR{(RADIUS).} בקצרה: אי אפשר סתם לחבר מחשב לשקע ולקבל גישה – המתג חוסם אותך עד שאתה מאמת את עצמך מול השרת. זהו הבסיס ל-\enLR{Wi-Fi} ארגוני \enLR{(WPA2-Enterprise)} ולבקרת גישה לרשת \enLR{(NAC)}, שנרחיב עליהם בפרק \enLR{6.}\par

\begin{storybox}\boxtitle{מהעולם האמיתי: למה ה-\enLR{Wi-Fi} באוניברסיטה שואל שם משתמש}\par  ברשת \enLR{eduroam} (שפועלת באוניברסיטאות בעולם, כולל בישראל) סטודנט מחיפה יכול להתחבר ל-\enLR{Wi-Fi} באוניברסיטה בברלין עם החשבון של הטכניון. איך? המחשב שולח שם משתמש לנקודת הגישה \enLR{(Authenticator)}, היא מעבירה את הבקשה לשרת \enLR{RADIUS} בברלין, שרואה את הסיומת "@\enLR{technion.ac.il"} ומעביר את השאלה לשרת \enLR{RADIUS} בטכניון, שמאמת ומחזיר "אשר". אף שרת לא צריך להכיר את כל המשתמשים בעולם – רק לדעת למי להעביר. זו בדיוק הסיבה ש-\enLR{RADIUS} הוא "פרוטוקול גישת המשתמשים", בעוד \enLR{TACACS}+ נשאר לניהול הציוד.\end{storybox}

\sectionheading{\enLR{3.8} בדיקה ופתרון תקלות}

אחרי שמגדירים, איך יודעים שהכול עובד – ואיך מאבחנים כשלא? הנה הכלים המרכזיים ותקלות נפוצות:\par

\begin{codebox}\begin{english}\codefont\footnotesize\color{cfg}\setlength{\parindent}{0pt}
R1\# \textcolor{cbold}{\textbf{show aaa sessions}}                       \textcolor{ccom}{\texthebrew{! מי מחובר כרגע דרך \enLR{AAA}}}\\
R1\# \textcolor{cbold}{\textbf{test aaa group tacacs+ admin P@ss legacy}}  \textcolor{ccom}{\texthebrew{! בודק אימות בלי להתחבר בפועל}}\\
R1\# \textcolor{cbold}{\textbf{debug aaa authentication}}                \textcolor{ccom}{\texthebrew{! מעקב חי אחר תהליך האימות}}
\end{english}\end{codebox}

\begin{itemize}
\item \textbf{הנתב "נתקע" כמה שניות ואז מקבל את המשתמש המקומי:} סימן שהשרת אינו נגיש (בדקו כתובת, \enLR{ACL}, ושהשירות פועל). הגיבוי ל-\code{local} עבד – זה תקין, אבל צריך לתקן את הקשר לשרת.
\item \textbf{"אימות נכשל" מיד, גם עם סיסמה נכונה:} לרוב המפתח המשותף \enLR{(key)} אינו זהה בין הנתב לשרת, או שהשרת לא מכיר את כתובת הנתב כלקוח מורשה.
\item \textbf{המשתמש נכנס אך לא מקבל הרשאת מנהל:} חסרה שורת ה-\code{aaa authorization exec}, או שהשרת לא מחזיר את רמת ההרשאה.
\end{itemize}

\sectionheading{\enLR{3.9} תרגילים – עם פתרונות}

\begin{exbox}\boxtitle{תרגיל \enLR{1} – איזה \enLR{A}?}\par  לכל משפט, ציינו אם מדובר באימות, הרשאה או רישום: (א) "המשתמש \enLR{dana} ניסתה להקליד \code{configure terminal} ונדחתה". (ב) "החשבון של \enLR{yossi} נעול אחרי \enLR{3} ניסיונות כושלים". (ג) "בדוח החודשי: \enLR{4,200} חיבורי \enLR{VPN}, ממוצע \enLR{38} דקות". \par\vskip2pt\hrule\vskip3pt\textbf{}(א) הרשאה (נבדק מה מותר לה לעשות). (ב) אימות (הוכחת זהות בכניסה). (ג) רישום (תיעוד פעולות).\end{exbox}

\begin{exbox}\boxtitle{תרגיל \enLR{2} – מה יקרה?}\par  נתב מוגדר: \code{aaa authentication login default group tacacs+ local}. (א) שרת ה-\enLR{TACACS}+ עונה "סיסמה שגויה". (ב) שרת ה-\enLR{TACACS}+ לא עונה כלל, ויש משתמש מקומי. מה קורה בכל מקרה? \par\vskip2pt\hrule\vskip3pt\textbf{}(א) המשתמש נדחה – אין מעבר ל-\enLR{local} אחרי דחייה מפורשת. (ב) אחרי \enLR{timeout} הנתב עובר ל-\enLR{local} ומאמת מולו – כי השרת "לא זמין".\end{exbox}

\begin{exbox}\boxtitle{תרגיל \enLR{3} – בחירת פרוטוקול}\par  לכל תרחיש בחרו \enLR{RADIUS} או \enLR{TACACS}+ ונמקו: (א) אימות \enLR{2,000} תלמידים ל-\enLR{Wi-Fi} של בית הספר. (ב) \enLR{5} מנהלי רשת שרוצים שמתמחה יוכל להקליד רק פקודות \code{show}. \par\vskip2pt\hrule\vskip3pt\textbf{}(א) \enLR{RADIUS} – גישת משתמשים לרשת \enLR{(802.1X).} (ב) \enLR{TACACS}+ – הוא היחיד שתומך באישור נפרד לכל פקודה.\end{exbox}

\begin{exbox}\boxtitle{תרגיל \enLR{4} – תקלה}\par  מנהל הקליד \code{aaa new-model} ואז \code{aaa authentication login default group tacacs+} (בלי \code{local}). השרת נפל, והמנהל התנתק. מה קרה, ואיך מונעים זאת? \par\vskip2pt\hrule\vskip3pt\textbf{}הוא ננעל מחוץ לנתב – אין שיטת גיבוי כשהשרת לא זמין. השחזור דורש גישה פיזית לקונסול והליך שחזור סיסמה. המניעה: תמיד לסיים את רשימת השיטות ב-\code{local} ולהגדיר משתמש מקומי מראש.\end{exbox}

\sectionheading{בנק שאלות שתלמידים שואלים – ותשובות מוכנות}

שאלות נפוצות בפרק ה-\enLR{AAA}, עם תשובות מוכנות.\par

\subheading{מושגי יסוד}

\textbf{ש: "אימות והרשאה נשמעים אותו דבר – מה ההבדל?"}\newline ת: משל שדה התעופה: \textbf{אימות} = בדיקת הדרכון (מי אתה). \textbf{הרשאה} = כרטיס הטיסה (לאיזה מטוס מותר לך לעלות). קודם מוודאים מי אתה, ורק אז קובעים מה מותר לך. אי אפשר הרשאה בלי אימות.\par

\textbf{ש: "ה-\enLR{key} בין הנתב לשרת – זו הסיסמה של המשתמש?"}\newline ת: לא! ה-\enLR{key} הוא סוד משותף בין ה\underline{נתב} ל\underline{שרת} (מצפין ומאמת את הקשר ביניהם). סיסמאות המשתמשים נמצאות בשרת. שני דברים שונים לגמרי.\par

\textbf{ש: "\enLR{MFA} ו-\enLR{2FA} זה אותו דבר?"}\newline ת: \enLR{2FA} (שני גורמים) הוא מקרה פרטי של \enLR{MFA} (רב-גורמי). העיקרון: שילוב גורמים מ\underline{סוגים שונים} – משהו שאתה יודע (סיסמה) + משהו שיש לך (טלפון) + משהו שאתה (טביעת אצבע). שתי סיסמאות אינן \enLR{MFA.}\par

\subheading{מקומי מול שרת, ותקלות}

\textbf{ש: "אם שרת ה-\enLR{AAA} נופל, ננעלנו החוצה מכל הנתבים?"}\newline ת: לא, אם הגדרתם נכון: רשימת השיטות מסתיימת ב-\code{local} (למשל \code{... group tacacs+ local}). כשהשרת לא זמין, הנתב עובר אוטומטית לאימות מול משתמש מקומי. לכן \underline{תמיד} משאירים גיבוי מקומי.\par

\textbf{ש: "מה קורה אם הקלדתי \code{aaa new-model} בטעות בלי משתמש מקומי?"}\newline ת: אתם עלולים להינעל החוצה – כי מרגע הפקודה כל הקווים עוברים ל-\enLR{AAA}, ואם אין משתמש ואין רשימת ברירת מחדל, אין כניסה. כלל ברזל: קודם \code{username}, אחר כך \code{aaa new-model}, ולהשאיר חלון פתוח עד שבודקים בחלון שני.\par

\textbf{ש: "אם השרת עונה 'סיסמה שגויה', הנתב עובר לאימות מקומי?"}\newline ת: לא! מעבר לשיטה הבאה קורה רק כשהשרת \underline{לא זמין} \enLR{(timeout).} דחייה מפורשת מהשרת = המשתמש נדחה, בלי מעבר. זו נקודה שנשאלת בבגרות.\par

\subheading{\enLR{RADIUS} מול \enLR{TACACS}+}

\textbf{ש: "אז איזה מהם 'יותר טוב', \enLR{RADIUS} או \enLR{TACACS}+?"}\newline ת: תלוי במשימה. ל\underline{ניהול ציוד} (מנהלים שמתחברים לנתבים, עם אישור לכל פקודה) – \enLR{TACACS}+ עדיף. ל\underline{גישת משתמשים לרשת} \enLR{(Wi-Fi/802.1X, VPN)} – \enLR{RADIUS.} בפועל ארגונים משתמשים בשניהם, כל אחד לתפקידו.\par

\textbf{ש: "למה \enLR{RADIUS} על \enLR{UDP} אם \enLR{TCP} אמין יותר?"}\newline ת: \enLR{RADIUS} תוכנן ב-\enLR{1991} לשרתי דיאל-אפ עם אלפי חיבורים; \enLR{UDP} קל ומהיר, והלקוח פשוט שולח שוב אם אין תשובה. \enLR{TACACS}+ בחר \enLR{TCP} כי הוא מנהל שיחה רב-שלבית ורוצה לדעת מיד אם השרת נפל. שניהם עובדים היום.\par

\textbf{ש: "אומרים ש-\enLR{RADIUS '}לא עושה הרשאה' – אבל הוא כן מחזיר הרשאות, לא?"}\newline ת: נכון, זו נקודה עדינה. \enLR{RADIUS} כן נושא \enLR{attributes} של הרשאה \enLR{(VLAN}, רמה) בתוך תשובת ה-\enLR{Accept}, אבל הוא \underline{משלב} אימות והרשאה בצעד אחד ואינו תומך באישור \underline{נפרד לכל פקודה} כמו \enLR{TACACS+.} בבגרות "אינו מתייחס לניהול הרשאות" מתקבל כנכון במובן הזה.\par

\textbf{ש: "מה זה \enLR{802.1X} במשפט אחד?"}\newline ת: תקן שבו המתג/נקודת הגישה \underline{חוסמים את הפורט} עד שהמכשיר מאמת את עצמו מול שרת \enLR{RADIUS.} זה מה שמונע ממישהו לחבר מחשב לשקע ולקבל גישה, וזה הבסיס ל-\enLR{Wi-Fi} ארגוני \enLR{(WPA2-Enterprise).}\par

\textbf{ש: "אפשר לחבר \enLR{AAA} ל-\enLR{Active Directory} של בית הספר?"}\newline ת: כן – זה בדיוק מה שעושים בארגונים. שרת ה-\enLR{AAA (ISE/NPS/FreeRADIUS)} מקבל את הבקשה מהנתב ומאמת מול ה-\enLR{AD/LDAP.} כך המורים משתמשים באותם שם וסיסמה בכל מקום.\par

\sectionheading{מילון מונחים – הגדרות}

הגדרות תמציתיות של המונחים המרכזיים בפרק. שימושי גם כדף עזר לבחינה (מותר מילון מונחים).\par

\begin{xltabular}{\linewidth}{|>{\setRTL\raggedleft\arraybackslash}X|>{\setRTL\raggedleft\arraybackslash}X|}
\hline
\rowcolor{hdrbg}\textbf{הגדרה} & \textbf{מונח} \\
\hline
\endhead
מסגרת של שלושה שירותים: \enLR{Authentication, Authorization, Accounting.} & \textbf{\enLR{AAA}} \\
\hline
קביעת זהות המשתמש – מי אתה. & \textbf{\enLR{Authentication} (אימות)} \\
\hline
קביעת מה מותר למשתמש לעשות אחרי שאומת. & \textbf{\enLR{Authorization} (הרשאה)} \\
\hline
תיעוד פעולות המשתמש: מתי נכנס, אילו פקודות הקליד. & \textbf{\enLR{Accounting} (רישום)} \\
\hline
אימות מול מסד משתמשים בקונפיג של הנתב עצמו. & \textbf{\enLR{AAA} מקומי} \\
\hline
אימות מול שרת מרכזי \enLR{(ISE/ACS)} בפרוטוקול \enLR{RADIUS} או \enLR{TACACS+.} & \textbf{\enLR{AAA} מבוסס-שרת} \\
\hline
ההתקן (נתב/מתג) שפונה לשרת ה-\enLR{AAA} בשם המשתמש. & \textbf{\enLR{NAS / AAA client}} \\
\hline
פרוטוקול \enLR{AAA} פתוח, \enLR{UDP (1812/1813)}; מצפין רק את הסיסמה; לגישת משתמשים. & \textbf{\enLR{RADIUS}} \\
\hline
פרוטוקול \enLR{AAA} של סיסקו, \enLR{TCP 49}; מצפין את כל החבילה; מפריד שירותים; לניהול ציוד. & \textbf{\enLR{TACACS}+} \\
\hline
סוד משותף בין הנתב לשרת המצפין ומאמת את הקשר ביניהם (לא סיסמת המשתמש). & \textbf{\enLR{key / shared secret}} \\
\hline
שרתי ה-\enLR{AAA} של סיסקו \enLR{(ACS} ישן, \enLR{ISE} מודרני). & \textbf{\enLR{ACS / ISE}} \\
\hline
הפקודה שמפעילה את מסגרת ה-\enLR{AAA} בנתב. & \textbf{\enLR{aaa new-model}} \\
\hline
רשימת שיטות אימות מסודרת; עוברים לשיטה הבאה רק אם הקודמת לא זמינה. & \textbf{\enLR{method list}} \\
\hline
תקן שבו המתג/\enLR{AP} חוסם פורט עד לאימות מול שרת \enLR{RADIUS.} & \textbf{\enLR{802.1X}} \\
\hline
ב-\enLR{802.1X}: המכשיר המבקש גישה \enLR{(supplicant)} והמתג שחוסם \enLR{(authenticator).} & \textbf{\enLR{Supplicant / Authenticator}} \\
\hline
אימות רב-שלבי – שילוב גורמים מסוגים שונים (יודע/יש/הוא). & \textbf{\enLR{MFA / 2FA}} \\
\hline
\end{xltabular}
\end{document}
